\chapter{Introduction}
\label{chap:intro}

Complexitylib is a Lean 4 formalization of computational complexity theory,
built on Mathlib. This blueprint is its map. Each chapter covers one area; each
node is a definition or result that is either formalized, with a link to its
Lean declaration, or planned, with the results it depends on. The dependency
graph shows which planned results are ready to be worked on.

The library's machine model follows Arora and Barak: multi-tape Turing machines
over the alphabet $\{0, 1, \sqcup, \triangleright\}$ with a read-only input
tape, work tapes, and an output tape. Complexity classes are defined from this
model, so every class-level theorem is a statement about concrete machines.
Other models (Boolean circuits, random access machines, interactive protocols,
logics) are related to it by explicit simulations.

\section{Reading the blueprint}

\begin{itemize}
  \item A node with a Lean link is formalized. Its prose is meant to state the
    linked Lean statement with all of its hypotheses, but the Lean statement is
    authoritative: where the two differ, what is proved is the Lean statement.
    CI checks that every declaration of the library depends on no axioms
    beyond Lean's three standard ones, so no declaration uses \texttt{sorry}
    or a custom axiom. It also checks the blueprint's links and markers: every
    Lean name exists; a statement is marked formalized exactly when it has a
    Lean link; the proof of a linked result is marked formalized; a theorem
    cites at least one Lean proof, not only definitions; and every label,
    \emph{uses} entry, and cross-reference resolves. Nothing checks that the
    prose matches the Lean statement; that is done by review.
  \item A node without a Lean link is planned. Its \emph{uses} list names the
    results a proof is expected to need.
  \item A result proved under a hypothesis is stated as that implication. For
    example, the library proves that $\cc{NL} \subseteq \cc{coNL}$ implies
    $\cc{NL} = \cc{coNL}$; the unconditional Immerman--Szelepcs\'enyi theorem is
    a separate planned node. A proved implication never stands in for the
    theorem it would give.
  \item Statements fix their conventions: resource bounds are explicit,
    encodings are concrete, and malformed inputs are handled. When a textbook
    result has several variants, the node says which one is formalized.
\end{itemize}

\section{How results are proved}

Most class-level results in the library are best proved through a small number
of characterization theorems rather than by building a new machine:

\begin{itemize}
  \item \textbf{Cobham's theorem} (Theorem~\ref{thm:cobham}) makes $\cc{FP}$ a
    programming language: a function is polynomial-time exactly when it is
    generated by composition and bounded recursion on notation from a few base
    functions. Closure lemmas turn this into a toolkit for writing
    polynomial-time functions as ordinary Lean terms.
  \item \textbf{The witness characterization of $\cc{NP}$}: a language with a
    polynomial-time verifier and polynomially bounded witnesses is in
    $\cc{NP}$.
  \item \textbf{Iteration in polynomial space}: iterating a polynomial-time
    step function exponentially many times on a polynomial-size state stays in
    $\cc{PSPACE}$. Savitch's theorem and $\cc{IP} \subseteq \cc{PSPACE}$ are
    proved this way.
  \item \textbf{Universal simulation} (Theorem~\ref{thm:utm}): one fixed
    machine simulates every machine through the computable compiler
    $p \mapsto \langle \alpha, p \rangle$, preserving and reflecting halting and
    output, with polynomial overhead.
\end{itemize}

The main gap is logarithmic space, which has no comparable programming layer
yet; see the chapter on space complexity.

\section{Status and priorities}

Headline formalized results include Cook--Levin
(Theorem~\ref{thm:cook-levin}), Cobham's theorem, universal simulation, and the
deterministic time hierarchy (Theorem~\ref{thm:time-hierarchy}); each chapter
lists the rest. Infrastructure priorities that are not mathematical results,
such as proof automation, consolidating duplicated encodings, and CI checks for
unproved hypotheses, are described in \texttt{ROADMAP.md} in the repository.
